\documentclass[10pt]{article}
\usepackage[T1]{fontenc}
\usepackage[utf8]{inputenc}
\usepackage[french]{babel}
\usepackage{amsmath,amssymb,mathrsfs,mathpazo}
\usepackage{geometry}
\geometry{a4paper,margin=2cm}
\usepackage{graphicx,array,tabularx,tabularray,colortbl,multirow,multicol,booktabs,enumitem,pifont,ulem,fancybox,diagbox,cancel}
\usepackage{pstricks,pstricks-add,pst-plot,pst-tree,pst-node,pst-text,pst-func,pst-math,pst-bspline,pst-3dplot,pst-eucl,pst-solides3d}
\let\pscalculate\undefined
\usepackage{pst-bezier}
\usepackage{tikz,pgfplots,tkz-tab}
\usetikzlibrary{arrows,arrows.meta,patterns,calc,positioning,trees,decorations.markings,shapes.geometric,shapes.misc,shapes.arrows,intersections,decorations.pathmorphing,decorations.pathreplacing,backgrounds,shadows,babel}
\pgfplotsset{compat=1.18}
\usepgfplotslibrary{fillbetween}
\usepackage[np]{numprint}
\usepackage{listings,makecell,pas-tableur,twoopt}
\let\barre\undefined
\usepackage[pstricks]{tablvar}
\definecolor{darkgreen}{RGB}{0,100,0}
\providecommand{\red}{\color{red}}
\providecommand{\blue}{\color{blue}}
\providecommand{\green}{\color{green}}
\providecommand{\black}{\color{black}}
\providecommand{\white}{\color{white}}
\providecommand{\gray}{\color{gray}}
\providecommand{\cyan}{\color{cyan}}
\providecommand{\magenta}{\color{magenta}}
\providecommand{\yellow}{\color{yellow}}
\providecommand{\violet}{\color{violet}}
\providecommand{\orange}{\color{orange}}
\providecommand{\euro}{\texteuro}
\providecommand{\EUR}{\texteuro}
\providecommand{\ieme}{\textsuperscript{e}}
\providecommand{\degres}{\ensuremath{^\circ}}
\providecommand{\degree}{\ensuremath{^\circ}}
\providecommand{\vv}[1]{\overrightarrow{#1}}
\providecommand{\Coord}[1]{#1}
\providecommand{\decosix}{}
\providecommand{\starredbullet}{\ensuremath{\bullet}}
\pagestyle{empty}
\setlength{\parindent}{0pt}
\usepackage{adjustbox}\DeclareUnicodeCharacter{2236}{:}\DeclareUnicodeCharacter{2190}{\ensuremath{\leftarrow}}
\providecommand{\R}{\mathbb{R}}
\providecommand{\N}{\mathbb{N}}
\providecommand{\Z}{\mathbb{Z}}
\providecommand{\Q}{\mathbb{Q}}
\providecommand{\D}{\mathbb{D}}
\providecommand{\C}{\mathbb{C}}
\providecommand{\e}{{\mathrm{e}}}
\providecommand{\vect}[1]{\overrightarrow{#1}}
\providecommand{\vectt}[1]{\overrightarrow{\mathrm{#1}}}
\providecommand{\vvt}[1]{\overrightarrow{\mathrm{#1}}}
\providecommand{\np}[2][]{{#2~#1}}
\providecommand{\vv}[1]{{\overrightarrow{#1}}}
\renewcommand{\d}{{\mathrm{d}}}
\providecommand{\degres}{{^\circ}}
\providecommand{\degre}{{^\circ}}
\providecommand{\blue}{\color{blue}}
\providecommand{\red}{\color{red}}
\providecommand{\enskip}{\quad}
\providecommand{\llbracket}{\lbrack\!\lbrack}
\providecommand{\rrbracket}{\rbrack\!\rbrack}
\providecommand{\euro}{{\text{€}}}
\providecommand{\no}{{\text{n°}}}
\providecommand{\interval}[2]{{[#1;#2]}}
\providecommand{\intervFO}[2]{{[#1;#2[}}
\providecommand{\dotfill}{\ldots}

\providecommand{\ds}{\displaystyle}
\providecommand{\textding}[1]{#1}
\providecommand{\Oij}{\(\mathrm{O};\vec{\imath},\vec{\jmath}\)}
\providecommand{\Oijk}{\(\mathrm{O};\vec{\imath},\vec{\jmath},\vec{k}\)}
\providecommand{\Ouv}{\(\mathrm{O};\vec{u},\vec{v}\)}
\providecommand{\intervFF}[2]{\ensuremath{[#1;#2]}}
\renewcommand{\euro}{\ifmmode\text{\texteuro}\else\texteuro\fi}
\DeclareUnicodeCharacter{00B0}{\ifmmode^{\circ}\else\textdegree\fi}
\DeclareUnicodeCharacter{20AC}{\ifmmode\text{\texteuro}\else\texteuro\fi}
% Énoncés : Éducation nationale. Transcription et corrigés : APMEP, Denis Vergès et contributeurs. Découpage et classement : CoopMaths, Sébastien Lozano, Éric Elter et contributeurs. Adaptation de la présentation : Maths Explorer.
% CC BY-SA : https://coopmaths.fr/www/about/
% https://coopmaths.fr/alea/static/eam/2026/tex/eam_specialite_2026_06_asie_1_cor.tex
\begin{document}
\noindent\textit{APMEP / CoopMaths. CC BY-SA. Présentation adaptée par Maths Explorer.}\par\medskip
\subsection*{Partie A -- Étude d'un modèle discret}

	\begin{enumerate}
		\item $1\%$ de \np{1100} c'est $11$, donc on a {$u_1=\np{1100}-11=\np{1089}$.}

		Cela signifie que, selon ce modèle, le niveau du lac sera de \np[~m]{1089} au-dessus du niveau de la mer en 2021.

		\item
		\begin{enumerate}
			\item Diminuer une quantité de 1\% revient à multiplier celle-ci par \np{0.99}.

			Puisque pour tout entier $n$, le niveau à l'année $n+1$ est 1\% en dessous de celui à l'année $n$, la relation de récurrence de la suite {$(u_n)$} est :\quad {$\forall n \in \N, u_{n+1}=\np{0.99}u_n$}.

			\item La relation de récurrence est multiplicative, {$(u_n)$} est donc une suite géométrique, de raison \np{0.99}.

			\item On connaît l'expression explicite d'une suite géométrique :
			\quad {$\forall n \in \N, u_n=\np{1100}\times\np{0.99}^n$}.
		\end{enumerate}

		\item Selon ce modèle, le niveau atteint en 2028 sera égal à {$\np{1100}\times\np{0.99}^8$}.

		\emph{Remarque :} Avec une calculatrice, on arrive à environ \np[~m]{1015} : ce lac risque d'être complètement asséché, puisque son niveau aura diminué de 85~m, ce qui est la profondeur maximale estimée d'un des plus grands lacs de retenue d'Europe, le barrage de Serre-Ponçon dans les Hautes-Alpes.

		\item Ici, on cherche donc l'indice $n$ du premier terme de la suite $\big(u_n\big)$ qui sera \og en dessous de \np[m]{1020}\fg{}.

		On va choisir d'interpréter \og en dessous de \np[m]{1020}\fg{} comme signifiant \og inférieur ou égal à \np[m]{1020} \fg{}.

		Le programme est initialisé avec les variables \texttt{u} et \texttt{n} contenant respectivement $u_0$ et $0$.

		Tant que le terme stocké dans la variable \texttt{u} n'est pas inférieur ou égal à \np{1020}, autrement dit, tant qu'il est strictement supérieur à \np{1020}, on va passer à l'indice suivant (c'est la ligne \texttt{n = n + 1}), et au terme suivant, en appliquant la relation de récurrence.

		Voici le programme complété :
\begin{lstlisting}[style=python,xleftmargin=4.5cm,xrightmargin=4.5cm,numbers=none]
u = 1100
n = 0
while u > 1020:
	n = n + 1
	u = u*0.99
print(n)
\end{lstlisting}
	\end{enumerate}

\subsection*{Partie B -- Étude d'un modèle continu}
	\begin{enumerate}
		\item 2020 correspond à l'année 0, on calcule $f(0)$ :

		{$f(0) = \np{1050}\,\e^{-0,01\times0} + 50=\np{1050}\,\e^{0} + 50=\np{1050}\times1+50=\np{1100}$}.

		Ce modèle donne un niveau calculé à \np[~m]{1100} en 2020, ce qui est cohérent avec l'observation.

		\item On sait que la dérivée d'une fonction de la forme \quad {$x \longmapsto \e^{ax}$},\quad  où $a$ est un réel est \quad {$x \longmapsto a\e^{ax}$}.

		On a donc, pour tout réel positif $t$ :\quad  {$f'(t)=\np{1050}\times(-\np{0.01})\times \e^{-\np{0.01}t}+0$}

		soit, après simplification,\quad {$f'(t)=-\np{10.5}\,\e^{-\np{0.01}t}$}.

		\item La fonction exponentielle est à valeurs strictement positives, donc pour tout $t$ de l'intervalle {$[0;+\infty[$}, \quad {$f'(t)<0$}.

		La fonction $f$ est donc strictement décroissante sur {$[0;+\infty[$.}

		\item La production est impossible si le niveau du lac s'abaisse en-dessous de \np[~m]{1020} d'altitude.

		C'est ce qui se produira, selon ce modèle, à partir de 2028 (année 8 sur le graphique).
	\end{enumerate}

\needspace{10\baselineskip}

\end{document}
